\documentclass[10pt,a4paper]{amsart}
\usepackage[margin=19mm]{geometry}
\usepackage{amsmath,amssymb,mathtools}
\usepackage[scr=rsfs,cal=euler]{mathalfa}
\usepackage{xfrac}
\usepackage[unicode,hidelinks,bookmarks=false]{hyperref}
\newcommand{\ie}{i.e.\ }
\newcommand{\cf}{cf.\ }
\newcommand{\resp}{respectively\ }
\newcommand{\Subsection}{\subsection}
\providecommand{\nobreakdash}{\nobreak\hbox{-}}
\setlength{\emergencystretch}{2em}
\multlinegap0pt
\usepackage{bm}
\usepackage{bbm}
\usepackage{dsfont}
\usepackage{xspace}
\usepackage{cancel}
\usepackage{mathtools}
\usepackage{amsthm}
\makeatletter
\@ifundefined{theorem}{\newtheorem{theorem}{Theorem}}{}
\@ifundefined{proposition}{\newtheorem{proposition}[theorem]{Proposition}}{}
\makeatother
\newcommand{\G}{\tilde{G}}
\newcommand{\R}{\mathbb{R}}
\newcommand{\p}{\tilde{p}}
\title[Non-Euclidean Crystallographic Rigidity]{Non-Euclidean Crystallographic Rigidity}
\author{Jack Esson, Eleftherios Kastis, Bernd Schulze}
\date{Source: arXiv:2510.08128v2 (CC BY 4.0)}
\begin{document}
\maketitle

\noindent\emph{Source and licence.} Non-Euclidean Crystallographic Rigidity. Version arXiv:2510.08128v2 (CC BY 4.0), distributed under the Creative Commons Attribution 4.0 International licence, as recorded in the arXiv licence metadata for this exact version and shown on the abstract page https://arxiv.org/abs/2510.08128v2. Full rights notice: licenses/arxiv-2510.08128-CC-BY-4.0.txt.\par\smallskip

\noindent\textit{Editorial scope.} Counted unit: the Proposition whose source label is \texttt{LqCsJoiningExtension} in the article (source0.tex lines 669--671, source label \texttt{LqCsJoiningExtension}), delivered together with the article's own complete original proof of it (source0.tex lines 672--692, one \texttt{proof} environment of 21 lines). No mathematics is written, repaired or re-derived: statement and proof are the article's own text line for line. Required context retained uncounted: none. Every definition, display and label of those lines is retained unchanged. Omitted from this package and not redistributed: 1--668, 693--1554. Nothing inside a retained excerpt is omitted or altered: every retained body is delivered exactly as the article's lines are written, so no omission receipt of any kind accompanies this package. Citations: the retained excerpts carry no citation command at all, so no bibliography entry of the article is transcribed and no reference is redistributed. Reference adaptation: none was needed and none was made; every reference inside the delivered unit resolves inside it, so no unresolved reference or citation is produced. Printed slips kept literal: no printed word, symbol, hypothesis or number of the counted statement or of its proof was altered. Layout and class: the article's own class is \texttt{article} with packages including geometry, amsmath, amsthm, amssymb, amsfonts, graphicx, babel, enumerate, listings, bm, float, hyperref, bbm, tikz; the wrapper is the lane's pinned \texttt{amsart} editorial wrapper at 10pt on A4, which restates the theorem-like environments the retained text needs and adds the article's own macro definitions that the retained text needs (\texttt{\textbackslash G}, \texttt{\textbackslash R}, \texttt{\textbackslash p}), with extra wrapper packages (bm, bbm, dsfont, xspace, cancel, mathtools). The delivered numbering is the unit's own. Assets: no retained span contains a figure environment, a graphics-inclusion command, a PDF-page inclusion, a table or a TikZ picture; no image file is copied into this package, the package declares no asset dependency and no image file is redistributed with it. Licence: version arXiv:2510.08128v2 of this article is distributed under the Creative Commons Attribution 4.0 International licence, as recorded in the arXiv licence metadata for this exact version and shown on the abstract page https://arxiv.org/abs/2510.08128v2. A full rights notice is delivered at licenses/arxiv-2510.08128-CC-BY-4.0.txt. Characters: every retained excerpt is pure ASCII, so the delivered engine, XeLaTeX with its default Latin Modern text font, reports no missing character.
    \begin{proposition}\label{LqCsJoiningExtension}
		Let $(\G_1,\p_1)$ and $(\G_2,\p_2)$ be $\mathcal{C}_s$-symmetric frameworks in $(\R^2,\|\cdot\|_q)$ that are minimally $\mathcal{C}_s$-symmetrically infinitesimally rigid. Let $(G_1,m_1)$ and $(G_2,m_2)$ be the underlying $\mathcal{C}_s$-gain graphs of $(\G_1,\p_1)$ and $(\G_2,\p_2)$ respectively. Let $(G',m')$ be formed by a gained edge-joining move that combines $(G_1,m')$ and $(G_2,m')$. Let $(\G',\p')$ be a $\mathcal{C}_s$-regular framework derived from $(G',m')$. Then $(\G',\p')$ is minimally $\mathcal{C}_s$-symmetrically infinitesimally rigid.
	\end{proposition}

	\begin{proof}
		Suppose that the edge-joining move gives $G'= G_1\oplus G_2=(V(G_1)\cup V(G_2),E(G_1)\cup E(G_2)\cup\{(v_1,v_2;m'(e))\})$, where $v_1\in V(G_1)$ and $v_2\in V(G_2)$. It is clear that this move preserves $(2,1)$-tightness, so it remains only to prove that it preserves row-independence of the orbit rigidity matrix. Since translating a framework does not change its rigidity properties, without loss of generality, we may assume that $\tilde{p}_1(v_1)\neq m'(e)\tilde{p}_2(v_2)$, and  we may further assume that $\p'|_{V(\G_1)}=\p_1$ and $\p'|_{V(\G_2)}=\p_2$. If we prove that a framework following these assumptions is $\mathcal{C}_s$-symmetrically infinitesimally rigid, then $\mathcal{C}_s$-symmetric infinitesimal rigidity follows for all $\mathcal{C}_s$-regular configurations of $\G$. Observe that the orbit rigidity matrix of $(\G',\p')$ takes the following form:
		\begin{align*}
			\mathbf{O}_q(\G',\p',\mathcal{C}_s)=
			\begin{bmatrix}
				\mathbf{O}_q(\G_1,\p'|_{V(\G_1)},\mathcal{C}_s)&0\\
				0&\mathbf{O}_q(\G_2,\p'|_{V(\G_2)},\mathcal{C}_s)\\
				\begin{matrix}
					(\p'(v_1)-m'(e)\p'(v_2))^{(q-1)}&0
				\end{matrix}
				&
				\begin{matrix}
					(\p'(v_2)-(m'(e))^{-1}\p'(v_1))^{(q-1)}&0
				\end{matrix}
			\end{bmatrix}
			.
		\end{align*}
		It is clear that the rows for $E(G_1)\cup E(G_2)$ are linearly independent, so it remains only to show that the row for the new edge $e$ is independent of the others.

        The rank of $\mathbf{O}_q(\G_1,\p'|_{V(\G_1)},\mathcal{C}_s)\oplus\mathbf{O}_q(\G_2,\p'|_{V(\G_2)},\mathcal{C}_s)$ is $2|V(G)|-2$, so the resulting kernel is $2$-dimensional. One element of this kernel is the vector $u\in \R^{2|V(G)|}$ defined by setting $u(v)=(1,0)$ if $v\in V(G_1)$ and $u(v)=(0,0)$ if $v\in V(G_2)$. By regularity, it can be assumed that $\p'(v_1)-m'(e)\p'(v_2)$ has a non-zero horizontal component. From this, it follows that $e$ eliminates $u$ from the kernel, which leaves only a $1$-dimensional kernel, as required.
	\end{proof}

\end{document}
